
\documentclass[11pt]{article}
\usepackage[margin=1in]{geometry}
\usepackage[T1]{fontenc}
\usepackage{hyperref}
\usepackage{listings}
\usepackage{longtable}
\usepackage{enumitem}

\title{ROS 2 Architecture and Development Guide}
\author{}
\date{\today}

\begin{document}

\maketitle
\tableofcontents
\newpage

\section{Introduction}

ROS 2 (Robot Operating System 2) is a middleware framework for building distributed robotic systems.
ROS 2 is not an operating system. It is a collection of libraries, tools, communication middleware,
build systems, and conventions that simplify robot software development.

Key goals of ROS 2:

\begin{itemize}
\item Distributed communication
\item Real-time support
\item Security support
\item Multi-platform operation
\item Scalable architecture
\item DDS-based communication
\end{itemize}

\section{High-Level ROS 2 Architecture}

A ROS 2 application consists of:

\begin{itemize}
\item Nodes
\item Topics
\item Publishers
\item Subscribers
\item Services
\item Actions
\item Parameters
\item DDS Middleware
\end{itemize}

Architecture Flow:

\begin{verbatim}
+-------------+
| Application |
+-------------+
       |
       v
+-------------+
| rclcpp/rclpy|
+-------------+
       |
       v
+-------------+
| rcl Layer   |
+-------------+
       |
       v
+-------------+
| rmw Layer   |
+-------------+
       |
       v
+-------------+
| DDS Vendor  |
+-------------+
       |
       v
+-------------+
| Network     |
+-------------+
\end{verbatim}

\section{Core Concepts}

\subsection{Node}

A node is an executable process participating in a ROS graph.

Examples:

\begin{itemize}
\item Camera node
\item LiDAR node
\item Navigation node
\item Motor controller node
\end{itemize}

\subsection{Topic}

Topics provide asynchronous communication.

Example:

\begin{verbatim}
/camera/image
/scan
/cmd_vel
\end{verbatim}

\subsection{Publisher}

Produces messages.

\subsection{Subscriber}

Consumes messages.

\subsection{Service}

Synchronous request-response communication.

\subsection{Action}

Long-running operations with feedback.

\section{ROS 2 Middleware Stack}

\subsection{Application Layer}

User code written using:

\begin{itemize}
\item rclcpp (C++)
\item rclpy (Python)
\item rclc (C)
\end{itemize}

\subsection{RCL}

ROS Client Library abstraction.

Responsibilities:

\begin{itemize}
\item Node creation
\item Topic creation
\item Timer management
\item Executor interaction
\end{itemize}

\subsection{RMW}

ROS Middleware abstraction layer.

Responsibilities:

\begin{itemize}
\item DDS interaction
\item Serialization
\item Discovery
\end{itemize}

\subsection{DDS}

Data Distribution Service performs:

\begin{itemize}
\item Discovery
\item Transport
\item Reliability
\item QoS management
\end{itemize}

\section{Development in C}

The C client library is called rclc.

Typical usage:

\begin{enumerate}
\item Initialize context
\item Create node
\item Create publisher/subscriber
\item Create executor
\item Spin executor
\end{enumerate}

Advantages:

\begin{itemize}
\item Small memory footprint
\item Embedded support
\item Micro-ROS compatibility
\end{itemize}

\section{Development in C++}

The C++ client library is rclcpp.

Typical structure:

\begin{verbatim}
class MyNode : public rclcpp::Node
{
public:
    MyNode() : Node("my_node")
    {
    }
};
\end{verbatim}

Advantages:

\begin{itemize}
\item Strong typing
\item High performance
\item Production robot deployments
\end{itemize}

\section{Development in Python}

The Python client library is rclpy.

Typical structure:

\begin{verbatim}
class MyNode(Node):
    def __init__(self):
        super().__init__("my_node")
\end{verbatim}

Advantages:

\begin{itemize}
\item Rapid development
\item Easy prototyping
\item AI and ML integration
\end{itemize}

\section{ROS 2 Build System}

\subsection{ament}

ROS 2 uses ament as the build framework.

\subsection{colcon}

Build command:

\begin{verbatim}
colcon build
\end{verbatim}

Workspace structure:

\begin{verbatim}
workspace/
 ├── src/
 ├── build/
 ├── install/
 └── log/
\end{verbatim}

\section{Talker and Listener Example}

Example:

\begin{verbatim}
ros2 run demo_nodes_cpp talker
ros2 run demo_nodes_py listener
\end{verbatim}

\section{What Happens When Talker Starts?}

Step 1:

Operating system creates a process.

\begin{verbatim}
PID 1001
talker process
\end{verbatim}

Step 2:

main() executes.

Step 3:

rclcpp::init() initializes ROS.

Step 4:

Node object is created.

\begin{verbatim}
Node name = talker
\end{verbatim}

Step 5:

Publisher object is created.

\begin{verbatim}
Topic:
/chatter
\end{verbatim}

Step 6:

DDS participant is created.

Step 7:

DDS announces:

\begin{verbatim}
I publish topic /chatter
Type = std_msgs/String
\end{verbatim}

\section{What Happens When Listener Starts?}

Step 1:

Operating system creates listener process.

Step 2:

ROS initialization occurs.

Step 3:

Node object is created.

\begin{verbatim}
listener
\end{verbatim}

Step 4:

Subscriber created.

\begin{verbatim}
Topic:
/chatter
\end{verbatim}

Step 5:

DDS discovery starts.

Step 6:

Publisher and subscriber discover each other.

\section{Discovery Mechanism}

DDS performs automatic discovery.

Publisher broadcasts:

\begin{verbatim}
Topic Name
Message Type
QoS Profile
\end{verbatim}

Subscriber receives metadata.

Matching occurs.

Connection is established automatically.

\section{Message Flow}

Message path:

\begin{verbatim}
Talker Node
     |
Publisher
     |
rclcpp
     |
rcl
     |
rmw
     |
DDS Writer
     |
Network
     |
DDS Reader
     |
rmw
     |
rclpy
     |
Subscriber Callback
     |
Listener Node
\end{verbatim}

\section{Serialization}

Before transmission:

\begin{enumerate}
\item Message object created
\item Converted to binary stream
\item Passed to DDS
\item Sent over transport
\end{enumerate}

Example:

\begin{verbatim}
"Hello World: 10"
\end{verbatim}

becomes serialized bytes.

\section{Receiving a Message}

Listener side:

\begin{enumerate}
\item DDS receives bytes
\item DDS validates message
\item Deserialization occurs
\item Callback scheduled
\item Executor executes callback
\end{enumerate}

Example callback:

\begin{verbatim}
def callback(msg):
    print(msg.data)
\end{verbatim}

\section{Executors}

Executors run callbacks.

Single-threaded:

\begin{verbatim}
rclcpp::spin(node);
\end{verbatim}

Responsibilities:

\begin{itemize}
\item Wait for events
\item Process timers
\item Execute subscriptions
\item Execute services
\end{itemize}

\section{Quality of Service (QoS)}

QoS controls communication behavior.

Parameters:

\begin{itemize}
\item Reliability
\item Durability
\item History
\item Deadline
\item Lifespan
\end{itemize}

\section{Intra-Process Communication}

If nodes are in the same process:

\begin{itemize}
\item Copies can be reduced
\item Performance increases
\item Latency decreases
\end{itemize}

\section{Debugging Tools}

Useful commands:

\begin{verbatim}
ros2 node list
ros2 topic list
ros2 topic echo /chatter
ros2 topic info /chatter
ros2 doctor
\end{verbatim}

\section{Recommended Learning Path}

\begin{enumerate}
\item Learn nodes and topics
\item Learn publishers and subscribers
\item Learn services
\item Learn actions
\item Learn parameters
\item Learn launch files
\item Learn custom messages
\item Learn DDS and QoS
\item Learn lifecycle nodes
\item Build a complete robot application
\end{enumerate}

\section{Conclusion}

ROS 2 applications are collections of nodes communicating through DDS.
The talker/listener example demonstrates the complete communication chain:
process creation, node initialization, DDS discovery, message serialization,
transport, deserialization, callback execution, and application-level handling.

\end{document}
